\documentclass[11pt]{article}
\usepackage[margin=1in]{geometry}
\usepackage{graphicx}
\usepackage{amsmath}
\usepackage{amssymb}
\usepackage{booktabs}
\usepackage{natbib}
\usepackage{xcolor}
\usepackage{siunitx}
\usepackage{hyperref}
\usepackage{textcomp}
\usepackage[utf8]{inputenc}

\title{An artificial intelligence accelerated virtual screening platform
for drug discovery}
\author{}
\date{}

\begin{document}
\maketitle

\begin{abstract}
Structure-based virtual screening is a cornerstone of early drug discovery,
and the emergence of make-on-demand libraries containing billions of
purchasable compounds has dramatically expanded the chemical space available
for lead identification. Realising this opportunity, however, requires
docking programs that are simultaneously accurate enough to distinguish
true binders from close decoys and efficient enough to evaluate libraries
several orders of magnitude larger than those considered in the previous
generation of virtual screens. Most leading physics-based programs are
proprietary, free alternatives lag them on standard benchmarks, and recent
deep-learning scoring functions generalise poorly when the binding site is
known. Here we introduce RosettaGenFF-VS, an extension of the Rosetta
GALigandDock force field with additional atom types, new torsional
potentials, and an entropy term, together with the RosettaVS protocol that
couples a fast virtual-screening-express (VSX) mode to a flexible-receptor
virtual-screening-high-precision (VSH) mode for final re-ranking. We embed
these components in OpenVS, an open-source AI-accelerated virtual screening
platform that uses active-learning surrogates to triage multi-billion
compound libraries. RosettaGenFF-VS achieves a top 1\% enrichment factor
of $\mathrm{EF}_{1\%} = 16.72$ on CASF-2016, compared with $11.9$ for the
second-best physics-based method, and RosettaVS outperforms the second-best
method on DUD by a factor of two at 0.5\% and 1.0\% ROC enrichment. Applied
to two unrelated drug-discovery problems---the KLHDC2 ubiquitin ligase and
the VSD4 domain of human NaV1.7---each multi-billion compound campaign
completed in under seven days on a 3000-CPU HPC cluster with one RTX2080
GPU and yielded seven novel KLHDC2 hits (14\% hit rate) and four novel
NaV1.7 hits (44.4\% hit rate), all with single-digit micromolar binding
affinities, and a 2.0~\AA{} co-crystal structure of KLHDC2 with the best
initial hit validated the predicted binding pose.
\end{abstract}

\section{Introduction}

Structure-based virtual screening has become a central tool in early-stage
drug discovery because it can interrogate chemical space at a scale that
wet-lab assays alone cannot reach. The recent emergence of make-on-demand
libraries such as Enamine REAL and ZINC22, each cataloguing several billion
purchasable compounds, has expanded that accessible space by several
orders of magnitude relative to the historical high-throughput screening
decks of a few million molecules~\citep{lyu2019ultralarge,tingle2023zinc22,
irwin2005zinc}. Empirically, screens conducted at ultra-large scale often
recover more potent and more diverse chemotypes than screens of smaller
libraries, suggesting that the additional coverage translates into tangible
improvements in lead discovery~\citep{lyu2019ultralarge,gorgulla2020virtualflow}.
The challenge, however, is that the computational cost of evaluating
billions of poses with a physics-based docker scales with library size, so
that only a handful of multi-billion compound campaigns have been reported
to date~\citep{gorgulla2020virtualflow,gentile2022deepdocking}.

Several complementary strategies have emerged to bring ultra-large
screening within reach of modern high-performance computing clusters.
Parallelised docking platforms distribute independent calculations across
thousands of cores to absorb the wall-clock penalty of exhaustive
docking~\citep{gorgulla2020virtualflow,yaacoub2022ddgui}. Active-learning
and Bayesian-optimisation pipelines instead train a cheap surrogate on a
small fraction of docked compounds and use it to triage the remainder of
the library~\citep{graff2021accelerating,gentile2022deepdocking}.
Hierarchical screens apply successively more expensive scoring stages to
progressively smaller compound pools, and GPU-accelerated docking
engines compress wall-clock time further still. These approaches are
complementary, and each has moved the frontier of tractable screen size,
but all of them rest on the same underlying docking program; the accuracy
of that program sets a hard ceiling on the hit rate that even the most
efficient triage can achieve.

The accuracy question has, in turn, become a bottleneck for open science.
Leading proprietary docking programs such as Glide~\citep{friesner2004glide}
and GPU-accelerated variants integrated into community platforms like
VirtualFlow~\citep{gorgulla2020virtualflowgwo} offer strong benchmark
performance, but their licensing terms restrict unrestricted academic use
and preclude integration into fully open pipelines. AutoDock
Vina~\citep{trott2010vina} is widely adopted and freely available, yet it
trails commercial programs on standard scoring benchmarks. Deep-learning
scoring functions report competitive numbers on benchmarks such as
CASF-2016~\citep{su2019casf2016,liu2017pdbbind,scantlebury2023pointvs},
but their generalisation to novel compounds and previously unseen targets
is limited, and benchmark contamination remains a concern even when
Tanimoto similarity cut-offs of 0.6 on ligands and 30\% sequence identity
cut-offs on proteins are applied. When the binding site is known---as it
almost always is in a targeted virtual screen---physics-based methods
therefore remain the pragmatic choice for top-of-list
ranking~\citep{chaput2017usc,ebejer2019ligity,li2024cache}.

A practical consequence of this landscape is that there is no freely
available, scalable platform that combines an accurate physics-based force
field, explicit modelling of receptor flexibility, and active-learning
triage in a single open-source workflow. The community has tools that
excel in one or two of these dimensions, but the combination that a
multi-billion compound screen against a challenging target actually
requires has been missing. We argue that closing this gap calls for three
coordinated advances: a force field tuned specifically for
rank-ordering ligands against a shared target, a docking protocol that
offers both a fast screening mode and a flexible-receptor refinement
mode, and a surrounding platform that marshals both through active-learning
triage at scale.

In this work we describe such a combination and demonstrate it end-to-end
on two orthogonal drug-discovery problems. We introduce RosettaGenFF-VS,
an extension of the physics-based Rosetta generalised force
field~\citep{leman2020rosetta} that adds new atom types, new torsional
potentials, and an entropy term, so that the scoring function can
distinguish structurally distinct ligands competing for the same pocket.
Built on top of this force field, the RosettaVS protocol exposes two
docking modes: the virtual-screening-express (VSX) mode for fast initial
triage and the virtual-screening-high-precision (VSH) mode that restores
full receptor sidechain flexibility for final re-ranking of top-scoring
candidates. We wrap both modes inside OpenVS, an open-source
AI-accelerated virtual screening platform that trains a target-specific
neural-network surrogate online during the docking campaign, so that the
docker is only ever applied to the most promising fraction of a
multi-billion compound library. We benchmark the pipeline on CASF-2016
and the Directory of Useful Decoys, and then apply it to two unrelated
protein targets: the KLHDC2 ubiquitin ligase, which is emerging as a
PROTAC platform~\citep{slusarz2023klhdc3}, and the VSD4 domain of the
human voltage-gated sodium channel NaV1.7, a long-standing but
notoriously selectivity-limited analgesic
target~\citep{wang2015nav17methyl}. Against both targets we recover novel
single-digit micromolar hits in under seven days of compute, and a 2.0~\AA{}
X-ray co-crystal structure of KLHDC2 with our best initial hit validates
the predicted binding pose. Together these results establish that open,
physics-based, flexible-receptor virtual screening at billion-compound
scale is now tractable.

\section{Background and related work}

\paragraph{Ultra-large library docking.}
Make-on-demand libraries have expanded chemical space from the low millions
of compounds available in legacy screening decks to the several billions
of synthesisable molecules now catalogued by providers such as Enamine REAL
and ZINC22~\citep{lyu2019ultralarge,tingle2023zinc22}. Early ultra-large
screens demonstrated that both hit rates and chemotype diversity tend to
improve as the library grows, which motivated the development of
platforms able to evaluate billions of molecules
routinely~\citep{lyu2019ultralarge,gorgulla2020virtualflow}. VirtualFlow
illustrates this style of distributed ultra-large screening against
targets such as KEAP1~\citep{gorgulla2020virtualflow}, and later
extensions graft grey-wolf-optimised receptor flexibility onto the same
pipeline~\citep{gorgulla2020virtualflowgwo}. These efforts show that
ultra-large screening is possible in principle; our contribution is to
make it also open, accurate, and flexible-receptor.

\paragraph{Active-learning and deep-learning-guided triage.}
Rather than docking an entire library, active-learning approaches train a
cheap surrogate model on a small fraction of docked compounds and use it
to prioritise the rest. Pool-based Bayesian optimisation can recover a
substantial fraction of the top-ranked molecules in a 100-million-compound
library after docking only a few percent of
it~\citep{graff2021accelerating}. Deep Docking and its interface successors
deploy iterative neural-network surrogates in conjunction with a
conventional docker to accelerate ligand
ranking~\citep{gentile2022deepdocking,yaacoub2022ddgui}. Community
benchmarks such as CACHE suggest that while this acceleration is
reliable, hit rates on truly novel pockets remain
modest~\citep{li2024cache}, reinforcing the observation that the identity
and accuracy of the underlying scoring function dominate the success of
any active-learning wrapper.

\paragraph{Docking programs and scoring functions.}
Physics-based docking programs differ in their sampling strategies and in
how they combine force-field and empirical terms. Glide uses hierarchical
conformational search with an empirical-plus-force-field scoring function
and reports strong redocking accuracy on standard
benchmarks~\citep{friesner2004glide}. AutoDock Vina introduced a smoothed
empirical scoring function with multithreaded optimisation and remains a
widely used open-source reference point~\citep{trott2010vina}. The PDBbind
and CASF benchmarks evaluate scoring, ranking, docking, and screening
power on diverse protein-ligand
complexes~\citep{liu2017pdbbind,su2019casf2016}, and the Directory of
Useful Decoys assesses virtual-screening discrimination against
property-matched decoys~\citep{mysinger2012dude}. Comparative studies of
consensus docking and knowledge-based ligand methods have mapped both
the strengths and the characteristic failure modes of current
approaches~\citep{chaput2017usc,ebejer2019ligity}. Our RosettaVS
benchmarks plug directly into this ecosystem so that its performance can
be compared on equal terms with these reference methods.

\paragraph{Machine-learning scoring and structural limitations.}
Neural-network scoring functions have reported competitive CASF numbers,
yet concerns about train--test leakage persist even when ligand and
protein-similarity cut-offs are applied, and recent analyses argue that
careful dataset construction and interpretability testing are prerequisites
for treating these methods as drop-in replacements for physics-based
scoring~\citep{scantlebury2023pointvs}. A closely related issue is that
docking against apo protein structures typically degrades
virtual-screening performance because the pocket sidechains have not
rearranged to accommodate the ligand; this motivates flexible-receptor
refinement and holo-like pocket generation in the
docking pipeline~\citep{guterres2020apo}. RosettaVS addresses both
issues in a single workflow by staging a fast rigid-receptor screen
followed by a flexible-receptor refinement.

\paragraph{KLHDC2 and NaV1.7 as test targets.}
We chose KLHDC2 and NaV1.7 VSD4 precisely because they differ in almost
every biological and biophysical dimension. KLHDC2 is a substrate receptor
of the CUL2/CRL2 E3 ubiquitin ligase complex and recognises C-terminal
diglycine degrons through a compact propeller-domain
pocket~\citep{slusarz2023klhdc3}, which has made it an emerging platform
for PROTAC-style targeted protein degradation. The human voltage-gated
sodium channel NaV1.7, by contrast, is a large membrane protein whose
fourth voltage-sensing domain hosts a pocket for small-molecule
antagonists that stabilise an inactivated state of the channel, and where
subtle structural modifications of known aryl-sulfonamide scaffolds can
introduce new modes of gating
modulation~\citep{wang2015nav17methyl}. Together the two targets span a
soluble enzyme-like pocket and a membrane-embedded state-dependent site,
and therefore offer a stringent cross-check on the generality of the
platform.

\paragraph{Positioning of this work.}
Against this backdrop, RosettaGenFF-VS, RosettaVS, and OpenVS combine an
open, physics-based scoring function with flexible-receptor docking and
active-learning triage in a single workflow. Unlike deep-learning scoring
functions, the pipeline does not rely on the training-ligand distribution;
unlike proprietary commercial pipelines, all components are freely
available; and unlike previous active-learning approaches, receptor
flexibility is explicitly modelled at the final VSH stage so that induced-fit
effects can be captured in the scoring of the most promising candidates.
The benchmark and experimental sections that follow justify each of these
design choices in turn.

\section{Methods}

\subsection{RosettaGenFF-VS force field}
RosettaGenFF-VS is a virtual-screening-oriented extension of the Rosetta
generalised force field (RosettaGenFF) implemented within the Rosetta
GALigandDock framework~\citep{leman2020rosetta}. Starting from the existing
enthalpy formulation ($\Delta H$), which has previously delivered strong
docking accuracy on hundreds of thousands of complexes, we made three
additions tailored to ultra-large screening. First, the atom typing was
extended to cover chemical motifs commonly encountered in modern
synthesisable libraries; second, the torsional potential library was
expanded to reduce systematic errors on ligands with rotatable bonds
poorly represented in the original training set; and third, an entropy
estimator ($\Delta S$) was added so that the total free energy
$\Delta G = \Delta H - T\Delta S$ can rank-order structurally distinct
ligands competing for a shared target. Formally, the ranking score takes the
form
\begin{equation}
\Delta G \;=\; \Delta H \;-\; T\,\Delta S,
\end{equation}
where $\Delta H$ is the RosettaGenFF enthalpy term evaluated on the
docked complex and $\Delta S$ is the ligand-specific entropy estimator
introduced in RosettaGenFF-VS. An improved preprocessing script
complements these changes by automating parameter generation for large
libraries.

\subsection{RosettaVS protocol: VSX and VSH modes}
Building on RosettaGenFF-VS, the RosettaVS protocol defines two paired
docking modes. The VSX (virtual-screening express) mode is optimised for
rapid initial screening: it uses a rigid receptor and a compact sampling
schedule chosen to keep the per-ligand wall time compatible with
billion-compound triage. The VSH (virtual-screening high-precision) mode
is reserved for top candidates from VSX: it restores full receptor
sidechain flexibility, which is critical for pockets that undergo
induced conformational changes upon ligand binding. The two modes share
the same underlying force field, so that scores produced by VSX are
directly comparable with scores produced by VSH during re-ranking.

\subsection{OpenVS active-learning platform}
OpenVS is an open-source AI-accelerated virtual screening platform that
integrates RosettaVS with an active-learning surrogate. At each iteration,
a target-specific neural network is trained on the compounds that have
already been docked; the surrogate then ranks unseen compounds, and a
new batch of the highest-ranked molecules is sent to RosettaVS for
explicit docking. The new docking scores are appended to the training set
and the next iteration begins. This loop focuses expensive physics-based
docking on the narrow region of the library most likely to yield binders
while preserving the global coverage of the surrogate. The platform is
designed to be scalable and embarrassingly parallel, so that billion-scale
screens can be completed on modest HPC resources: both target campaigns
reported here ran on a cluster with 3000 CPUs and a single RTX2080 GPU.

\subsection{Benchmark protocols}
RosettaGenFF-VS was evaluated on the Comparative Assessment of Scoring
Functions benchmark (CASF-2016~\citep{su2019casf2016}), which contains 285
diverse protein-ligand complexes and provides precomputed decoy
conformations that decouple scoring from sampling. The docking-power test
measures the rate at which a scoring function ranks at least one pose with
RMSD $< 2$~\AA{} within the top one, two, or three decoys; the
screening-power test measures (i) the enrichment factor at a given cut-off
of ranked compounds, and (ii) the success rate of recovering the best
binder inside the top 1\%, 5\%, or 10\% of ranked ligands across all
targets. RosettaVS was additionally evaluated on the Directory of Useful
Decoys (DUD~\citep{mysinger2012dude}), which spans 40 pharmaceutically
relevant targets and more than 100{,}000 small molecules; performance is
reported as AUC of the receiver-operating-characteristic curve and as
ROC enrichment at 0.5\%, 1\%, and 2\% false-positive rates, the latter
being the regime most relevant to large-scale screening where only tens
to hundreds of compounds can be advanced to experimental testing.

\subsection{KLHDC2 virtual screening and experimental assays}
The KLHDC2 campaign used the OpenVS platform with VSX in RosettaVS to
screen the Enamine REAL library, which contains approximately 5.5 billion
purchasable compounds with a reported 80\% synthesis success rate, against
the target protein structure. The active-learning loop was run for ten
iterations, at which point approximately 6 million compounds---0.11\% of
the library---had been subjected to docking. The top-ranked 50{,}000
compounds from this VSX stage were re-docked with VSH to account for
induced conformational changes of pocket sidechains. The top 1{,}000
compounds from the VSH re-ranking were then filtered for predicted
solubility and for satisfied hydrogen bonds in the bound conformation,
and clustered to reduce redundancy; 54 representatives that passed the
filters and clustering were examined manually in PyMOL for favourable
geometries, 37 compounds were chosen for chemical synthesis, and 29 were
successfully synthesised. Activity was measured in an AlphaLISA
competition assay, in which each compound was tested for its ability to
displace a diglycine-containing SelK C-end degron peptide from KLHDC2,
and BioLayer Interferometry was used as an orthogonal confirmation.
Compound~29 (C29) was soaked into apo KLHDC2 crystals, and the complex
structure was solved to 2.0~\AA{} resolution. A focused follow-up screen
performed a substructure search on the acetyl-amino-methyl-triazole-acetic
acid backbone of C29 against ZINC22, which returned $\sim$381{,}567
compounds; of these, 21 were synthesised and re-tested in the AlphaLISA
assay with C29 as a positive control.

\subsection{NaV1.7 VSD4 virtual screening and electrophysiology}
The NaV1.7 campaign screened the ZINC22 library
($\sim$4.1 billion compounds) against the VSD4 domain of human NaV1.7
using the same VSX-then-VSH protocol. The active-learning loop was stopped
after the seventh iteration when the top predicted binding affinities had
converged. Approximately 4.5 million compounds ($0.11\%$ of the library)
were subjected to docking; the top 100{,}000 compounds were re-docked
with VSH. The re-ranked list was clustered, and the top 1{,}000
cluster representatives were filtered as in the KLHDC2 campaign; 160
compounds were examined manually and 10 were selected for synthesis after
excluding molecules containing the known aryl-sulfonamide warhead or
structurally resembling antihistamines and beta-blockers. Tanimoto
similarity of the selected molecules to known NaV1.7 inhibitors in ChEMBL
was less than 0.33, and 9 compounds were successfully synthesised.
Activity was measured in a whole-cell patch-clamp electrophysiology assay
on hNaV1.7 stably expressed in HEK-293 cells; concentration-response
curves were recorded at a holding potential of $-70$~mV for the
inactivated state and $-120$~mV for the resting state, and selectivity
was measured against hNaV1.5 (holding potential $-80$~mV) and the hERG
channel.

\section{Results}

\subsection{An open, AI-accelerated virtual screening platform}

Our first aim was to construct a physics-based virtual screening platform
that is simultaneously accurate, flexible in its treatment of the receptor,
and scalable to billion-compound libraries. The starting point is Rosetta
GALigandDock, which already supports flexible-receptor docking through the
RosettaGenFF force field but was originally tuned for screens of hundreds
of thousands of compounds rather than billions and lacks an entropy term
that would let it distinguish structurally distinct ligands competing for
the same pocket. We identified two deficiencies that must be corrected
before the approach can carry a multi-billion-compound screen: incomplete
coverage of functional groups in the atom-type and torsional tables, and
the absence of an entropy model in the scoring function.

To address both deficiencies we developed RosettaGenFF-VS, which extends
the RosettaGenFF atom typing and torsional potentials to cover functional
groups encountered in modern on-demand libraries, and adds an entropy
estimator ($\Delta S$) on top of the existing enthalpy model ($\Delta H$)
so that a full free-energy-like score can be used for ranking. We then
defined the RosettaVS protocol on top of this force field, which exposes
two coupled docking modes. VSX is optimised for throughput and is used
for the initial active-learning-guided sweep of the library; VSH restores
full receptor sidechain flexibility and is reserved for final re-ranking
of the top candidates from VSX. The two modes share the same underlying
force field, so that scores flow seamlessly from triage to refinement.

Even with these speed-ups, exhaustive docking of billions of compounds
remains prohibitive, so we wrapped RosettaVS in OpenVS, an open-source
platform that performs active-learning triage around the
docker~(Fig.~\ref{fig:pipeline}). At each iteration, a target-specific
neural network is trained on the compounds already docked and used to
rank the remainder of the library; the top-ranked unseen compounds are
then docked with RosettaVS, the new scores are appended to the training
set, and the cycle repeats. The platform is engineered to run in parallel
across a standard HPC cluster, and the two target campaigns reported
below each completed in under seven days on a 3000-CPU cluster with a
single RTX2080 GPU. The net effect of these three components---force
field, protocol, and platform---is a pipeline in which each layer is
tuned to the scale at which it is applied: coarse triage at
library scale, fast rigid-receptor docking at million scale, and
flexible-receptor refinement at ten-thousand scale. With the platform in
place, we next turned to quantitative benchmarks to confirm that the
accuracy of the underlying scoring function matches the best physics-based
methods in the community.

% Figure 5: OpenVS pipeline schematic
\begin{figure}[htbp]
\centering
\includegraphics[width=0.85\textwidth]{figures/fig_pipeline.pdf}
\caption{Schematic overview of the OpenVS platform pipeline. An ultra-large
compound library (billions of molecules) is triaged by an active-learning
neural-network surrogate, docked with VSX (fast mode), re-ranked with VSH
(flexible-receptor mode), filtered and clustered, then advanced to synthesis
and experimental assay.}
\label{fig:pipeline}
\end{figure}

\subsection{RosettaVS achieves state-of-the-art benchmark performance}

Having assembled the platform, we next asked whether RosettaGenFF-VS and
RosettaVS are competitive with the best physics-based methods on
community benchmarks. We focused on physics-based scoring functions because
recent deep-learning scoring functions, although they report strong
CASF-2016 numbers, have been shown to generalise poorly to novel compounds
and targets even after stringent similarity cut-offs (0.6 Tanimoto for
ligands and 30\% sequence identity for proteins) are applied.

We first ran the CASF-2016 docking-power and screening-power tests on
the 285-complex benchmark. On docking power, RosettaGenFF-VS achieves a
leading success rate at distinguishing the native binding pose from decoy
structures, and an accompanying binding-funnel analysis confirms that
the scoring potential drives conformational sampling towards the
lowest-energy minimum efficiently across a broad range of ligand RMSDs.
On screening power, RosettaGenFF-VS reaches a top 1\% enrichment factor
of $\mathrm{EF}_{1\%} = 16.72$, compared with $\mathrm{EF}_{1\%} = 11.9$
for the second-best physics-based method
(Fig.~\ref{fig:casf_ef1}, Table~\ref{tab:casf_screening}). It also
outperforms all other methods on the success rate of placing the best
binder inside the top 1\%, 5\%, and 10\% of ranked molecules, and
analysis by pocket-property subsets shows that the largest gains occur
on more polar, shallower, and smaller pockets---historically the
regimes where docking programs struggle.

% Figure 1: CASF-2016 screening power
\begin{figure}[htbp]
\centering
\includegraphics[width=0.55\textwidth]{figures/fig_casf_ef1.pdf}
\caption{CASF-2016 screening power comparison. Top 1\% enrichment factor
($\mathrm{EF}_{1\%}$) for RosettaGenFF-VS (16.72) versus the second-best
method (11.9), demonstrating a substantial margin in early-enrichment
performance on the 285-complex benchmark.}
\label{fig:casf_ef1}
\end{figure}

\begin{table}[htbp]
\centering
\caption{CASF-2016 screening power: top 1\% enrichment factor.}
\label{tab:casf_screening}
\begin{tabular}{lc}
\toprule
Method & $\mathrm{EF}_{1\%}$ \\
\midrule
RosettaGenFF-VS (this work) & 16.72 \\
Second-best method & 11.9 \\
\bottomrule
\end{tabular}
\end{table}

Together these benchmark results establish that RosettaGenFF-VS is a
leading scoring function for pose recognition and early enrichment on
fixed decoy sets. A practical virtual screen, however, also requires the
docker to sample good binding poses in the first place, which CASF
cannot test because its decoy conformations are precomputed; success on
CASF is therefore necessary but not sufficient to predict real-world
screening performance. To close this gap, we turned to the Directory of
Useful Decoys (DUD), which spans 40 pharmaceutically relevant targets
and more than 100{,}000 small molecules, and evaluated the full VSX and
VSH docking pipelines end-to-end. RosettaVS achieves the leading
AUC of the ROC curve among tested methods and outperforms the second-best
method by a factor of two at both 0.5\% and 1.0\% false-positive
ROC-enrichment rates, which are the cut-offs most relevant to ultra-large
screens where only tens to hundreds of compounds can be advanced to
experimental testing. VSH marginally outperforms VSX on DUD thanks to
its ability to model pocket sidechain rearrangements induced by the
ligand, confirming the value of staging a fast rigid-receptor sweep
against a flexible-receptor refinement.

\subsection{Discovery of novel KLHDC2 ubiquitin-ligase binders}

With the benchmarks established, we next asked whether the platform could
identify novel chemotypes for a target with no known drug-like
small-molecule binders. KLHDC2 is the substrate-receptor subunit of the
CUL2-RBX1 E3 ligase complex, it recognises C-terminal diglycine degrons
with nanomolar affinity through a KELCH-repeat propeller-domain pocket,
and its degron-binding site has recently emerged as a promising PROTAC
platform for targeted protein degradation. We therefore sought small
molecules that anchor to this diglycine-binding pocket.

We ran the OpenVS platform with VSX mode against the Enamine REAL library
($\sim$5.5 billion compounds, 80\% synthesis success rate) for ten
active-learning iterations. Across iterations the predicted binding
affinity of the top 0.1 percentile improved monotonically from
$-6.81$~kcal/mol in the first iteration to $-12.43$~kcal/mol in the tenth
(Fig.~\ref{fig:active_learning}), and no new global-minimum structures
appeared beyond the eighth iteration, so the screen was terminated at
iteration ten. The top 50{,}000 compounds were re-docked with VSH, and
the top 1{,}000 VSH compounds were filtered for predicted solubility and
hydrogen-bond completeness and clustered; 54 cluster representatives
passed the filters, of which 37 were selected for synthesis after manual
inspection of their predicted binding modes, and 29 were successfully
synthesised (Table~\ref{tab:klhdc2_campaign}). The entire campaign,
including all docking iterations and the VSH refinement, completed within
a week on a 3000-CPU HPC cluster equipped with a single RTX2080 GPU.

% Figure 2: Active-learning binding affinity convergence
\begin{figure}[htbp]
\centering
\includegraphics[width=0.6\textwidth]{figures/fig_active_learning.pdf}
\caption{Predicted binding affinity of the top 0.1 percentile compounds
across active-learning iterations for the KLHDC2 and NaV1.7 VSD4 campaigns.
KLHDC2: $-6.81$ kcal/mol (iteration~1) to $-12.43$ kcal/mol (iteration~10).
NaV1.7: $-10.8$ kcal/mol (iteration~1) to $-18.2$ kcal/mol (final iteration).}
\label{fig:active_learning}
\end{figure}

\begin{table}[htbp]
\centering
\caption{KLHDC2 virtual screening campaign summary.}
\label{tab:klhdc2_campaign}
\begin{tabular}{lc}
\toprule
Stage & Value \\
\midrule
Library (Enamine REAL) & $\sim 5.5 \times 10^9$ compounds \\
Compounds docked & $\sim 6 \times 10^6$ (0.11\%) \\
Active-learning iterations & 10 \\
Top 0.1 percentile iteration 1 affinity & $-6.81$~kcal/mol \\
Top 0.1 percentile iteration 10 affinity & $-12.43$~kcal/mol \\
VSH re-ranking depth & top 50{,}000 \\
Candidates after filter + cluster & 54 \\
Selected for synthesis & 37 \\
Synthesised & 29 \\
Initial hits (IC$_{50}$ low-\textmu M) & 7 \\
Focused screen substructure hits (ZINC22) & $\sim$381{,}567 \\
Focused screen synthesised & 21 \\
Additional focused-screen hits & 6 \\
Total KLHDC2 hits & 7 (14\% hit rate) \\
Wall-clock time & $<$~7~days on 3000 CPUs, 1 RTX2080 GPU \\
\bottomrule
\end{tabular}
\end{table}

Activity measurements in an AlphaLISA competition assay, in which each
compound was tested for its ability to displace a diglycine-containing
SelK C-end degron peptide from KLHDC2, identified compound~29 (C29) as
the best initial hit with an IC$_{50}$ of $\sim 3$~\textmu M, and the
same potency was subsequently confirmed by BioLayer Interferometry. To
reveal the binding mode of C29, we soaked apo KLHDC2 crystals with the
compound and determined the complex structure at 2.0~\AA{} resolution.
C29 binds the degron-binding pocket with its distal carboxyl group
contacting two critical arginine residues (Arg236 and Arg241) and a
serine residue (Ser269) that recognise the extreme C-terminus of the
native degron, its triazole moiety nestled among three aromatic residues
(Tyr163, Trp191, and Trp270) with a stabilising NH$\cdots$N hydrogen bond,
a hydrogen bond between Lys147 and the central carbonyl group, and its
tert-butylphenyl group packed into the auxiliary chamber of the pocket.
The central dimethyl-sulfide linker shows weaker electron density,
consistent with higher structural flexibility. The experimental binding
pose agrees closely with the computational prediction---superposition of
the experimental and docked structures shows the small molecule occupying
the same sub-pocket regions and making the same key polar contacts
predicted in silico---directly validating the accuracy of the RosettaVS
docking protocol on this target and confirming that the hit was
discovered via the correct binding mode rather than by coincidence.

To exploit the validated binding mode, we performed a focused follow-up
screen that ran a substructure search for the
acetyl-amino-methyl-triazole-acetic acid backbone of C29 against ZINC22,
which returned $\sim 381{,}567$ compounds. We docked these with the
GALigandDock flexible-docking mode, selected 21 of the top 100 compounds
after filtering, and synthesised and tested them in the AlphaLISA assay
using C29 as a positive control. Six of the 21 showed single-digit
micromolar IC$_{50}$ values, raising the total count of low-micromolar
KLHDC2 binders to seven and yielding an overall 14\% hit rate across
the combined initial and focused screens~(Fig.~\ref{fig:hit_rates}).
This result demonstrates that the platform can both discover novel
binders for a previously unliganded target and support structure-guided
expansion around a validated chemotype.

\subsection{Discovery of novel NaV1.7 VSD4 antagonists}

Having validated the platform on a soluble ubiquitin-ligase pocket, we
turned to the human voltage-gated sodium channel NaV1.7, which presents
a markedly different biophysical setting: a membrane-embedded channel
whose fourth voltage-sensing domain (VSD4) hosts a pocket for small
molecules that stabilise an inactivated state of the
channel~\citep{wang2015nav17methyl}. A pharmacologically useful
NaV1.7 VSD4 antagonist must not only bind but also exhibit state-dependent
inhibition and selectivity over related channels such as NaV1.5 and hERG,
which together raise a stringent bar for a virtual screen.

We ran the OpenVS platform with RosettaVS VSX mode against the ZINC22
library ($\sim 4.1 \times 10^9$ compounds) for seven active-learning
iterations, at which point the top predicted binding affinities had
converged. The predicted binding affinity of the top 0.1 percentile
improved from $-10.8$~kcal/mol in the first iteration to $-18.2$~kcal/mol
in the final iteration~(Fig.~\ref{fig:active_learning}). Approximately
4.5 million compounds---0.11\% of the library, mirroring the KLHDC2
campaign---were subjected to docking, and the top 100{,}000 were
re-ranked with VSH. We then clustered the VSH top list and applied the
same solubility and hydrogen-bond filters to the top 1{,}000 cluster
representatives; 160 molecules passed clustering and filtering and were
examined manually. To ensure novelty, we explicitly excluded molecules
containing the known aryl-sulfonamide warhead or structurally resembling
antihistamines and beta-blockers, so the 10 compounds ultimately
selected for synthesis all have Tanimoto similarity of less than 0.33 to
known NaV1.7 inhibitors in ChEMBL; 9 of the 10 were successfully
synthesised (Table~\ref{tab:nav17_campaign}).

\begin{table}[htbp]
\centering
\caption{NaV1.7 VSD4 virtual screening campaign summary.}
\label{tab:nav17_campaign}
\begin{tabular}{lc}
\toprule
Stage & Value \\
\midrule
Library (ZINC22) & $\sim 4.1 \times 10^9$ compounds \\
Compounds docked & $\sim 4.5 \times 10^6$ (0.11\%) \\
Active-learning iterations & 7 \\
Top 0.1 percentile iteration 1 affinity & $-10.8$~kcal/mol \\
Top 0.1 percentile final-iteration affinity & $-18.2$~kcal/mol \\
VSH re-ranking depth & top 100{,}000 \\
Manually examined candidates & 160 \\
Tanimoto similarity to known inhibitors & $<$~0.33 \\
Selected for synthesis & 10 \\
Synthesised & 9 \\
Hits (IC$_{50}$ $<$ 10~\textmu M) & 4 (44.4\% hit rate) \\
Best compound (Z8739902234) IC$_{50}$ & 1.32~\textmu M \\
\bottomrule
\end{tabular}
\end{table}

Activity was measured with whole-cell patch-clamp electrophysiology on
hNaV1.7 stably expressed in HEK-293 cells, using a $-70$~mV holding
potential for the inactivated state and a $-120$~mV holding potential for
the resting state. Four of the nine synthesised compounds showed
IC$_{50}$ values better than 10~\textmu M on the inactivated state,
corresponding to a 44.4\% hit
rate~(Fig.~\ref{fig:hit_rates}). The most potent compound,
Z8739902234, inhibited hNaV1.7 with an inactivated-state IC$_{50}$ of
1.32~\textmu M (95\% CI 1.14--1.55), and the second-best compound,
Z8739905023, showed an inactivated-state IC$_{50}$ of 2.23~\textmu M
(95\% CI 2.14--2.46)~(Fig.~\ref{fig:nav17_ic50},
Table~\ref{tab:nav17_ic50}).

% Figure 3: Hit rates
\begin{figure}[htbp]
\centering
\includegraphics[width=0.55\textwidth]{figures/fig_hit_rates.pdf}
\caption{Experimental hit rates for the KLHDC2 and NaV1.7 VSD4 virtual
screening campaigns. KLHDC2: 7 hits (14\% hit rate); NaV1.7: 4 hits with
IC$_{50} < 10$~\textmu M (44.4\% hit rate).}
\label{fig:hit_rates}
\end{figure}

Z8739902234 also displayed clear state dependence, with a resting-state
IC$_{50}$ of 13.1~\textmu M (95\% CI 12.78--13.45) compared with a
1.32~\textmu M inactivated-state IC$_{50}$ (95\% CI 1.14--1.55)---a
roughly ten-fold shift consistent with preferential binding to the
inactivated channel conformation. Selectivity profiling against two
clinically relevant off-target channels gave an hNaV1.5 IC$_{50}$ of
5.14~\textmu M (95\% CI 4.74--5.55) and an hERG IC$_{50}$ of 6.27~\textmu M
(95\% CI 5.27--7.37) for Z8739902234
at the hNaV1.7 reference IC$_{50}$ of 1.33~\textmu M (95\% CI 1.14--1.55),
giving a roughly four-fold window over NaV1.5 and a roughly five-fold
window over hERG~(Table~\ref{tab:nav17_ic50}). This selectivity profile
is modest but tractable, and the fact that it was obtained from a screen
whose selected compounds have Tanimoto similarity below 0.33 to known
NaV1.7 inhibitors underscores the platform's ability to access genuinely
novel chemotypes at a structurally subtle site.

% Figure 4: NaV1.7 IC50 values
\begin{figure}[htbp]
\centering
\includegraphics[width=0.7\textwidth]{figures/fig_nav17_ic50.pdf}
\caption{IC$_{50}$ values (mean, 95\% CI) for the two best NaV1.7 VSD4 lead
compounds. Z8739902234: inactivated-state IC$_{50}$ = 1.32~\textmu M
(1.14--1.55), resting-state IC$_{50}$ = 13.1~\textmu M (12.78--13.45);
selectivity against NaV1.5 (5.14~\textmu M, 4.74--5.55) and hERG
(6.27~\textmu M, 5.27--7.37). Z8739905023: inactivated-state IC$_{50}$ =
2.23~\textmu M (2.14--2.46).}
\label{fig:nav17_ic50}
\end{figure}

\begin{table}[htbp]
\centering
\caption{hNaV1.7 concentration-response and selectivity for lead compounds.
All values are IC$_{50}$ in \textmu M (mean, 95\% confidence interval).}
\label{tab:nav17_ic50}
\begin{tabular}{lcc}
\toprule
Compound / channel & IC$_{50}$ (\textmu M) & 95\% CI \\
\midrule
Z8739902234, hNaV1.7 inactivated ($-70$~mV) & 1.32 & 1.14--1.55 \\
Z8739902234, hNaV1.7 resting ($-120$~mV) & 13.1 & 12.78--13.45 \\
Z8739902234, hNaV1.7 reference & 1.33 & 1.14--1.55 \\
Z8739902234, hNaV1.5 ($-80$~mV) & 5.14 & 4.74--5.55 \\
Z8739902234, hERG & 6.27 & 5.27--7.37 \\
Z8739905023, hNaV1.7 inactivated & 2.23 & 2.14--2.46 \\
\bottomrule
\end{tabular}
\end{table}

\section{Discussion}

We have described a state-of-the-art, open-source, physics-based virtual
screening platform that combines a force field tuned for
rank-ordering competing ligands, a docking protocol that couples a fast
rigid-receptor mode to a flexible-receptor refinement, and an
active-learning triage layer that keeps billion-compound screens
tractable on modest HPC resources. Applied end-to-end, the platform
discovered seven novel low-micromolar binders to the KLHDC2 ubiquitin
ligase (14\% hit rate) and four novel single-digit-micromolar antagonists
to the VSD4 domain of human NaV1.7 (44.4\% hit rate), with the most
potent NaV1.7 compound inhibiting the inactivated channel at an
IC$_{50}$ of 1.32~\textmu M and a 2.0~\AA{} co-crystal structure
confirming the predicted KLHDC2 binding pose. Taken together, these
results show that open, accurate, flexible-receptor docking at
multi-billion-compound scale is now practical.

Several design choices in RosettaGenFF-VS and RosettaVS contribute to
the benchmark performance we observe. The combination of high docking
accuracy and efficient conformational sampling allows the scoring
function to find the correct binding minimum of a protein--ligand
complex more reliably than the physics-based baselines on CASF-2016:
the top 1\% enrichment factor of 16.72 is substantially larger than the
11.9 obtained by the next-best method, and the binding-funnel analysis
indicates that this is not a knife-edge result confined to a narrow RMSD
window. Equally important, the analysis by pocket subsets on the DUD
benchmark shows that the relative advantage of RosettaGenFF-VS is largest
on polar, shallower, and smaller pockets, which are historically where
docking programs are weakest. We attribute this behaviour to the
improved balance between protein--ligand and intra-ligand energies in
the extended force field, since these pocket classes rely less on
hydrophobic burial and more on specific polar contacts that must be
scored consistently across structurally diverse ligands.

The two target campaigns reinforce the benchmark results in a practical
setting, and at the same time illuminate where the underlying platform is
distinct from prior ultra-large screens. VirtualFlow and deep-docking
pipelines have both demonstrated that distributed rigid-receptor docking
or neural-network-guided triage can make billion-compound screens
tractable, but neither combines an open physics-based scoring function
with an explicit flexible-receptor refinement at the final
stage~\citep{gorgulla2020virtualflow,gentile2022deepdocking,
yaacoub2022ddgui}. Deep Docking and related platforms work well on
pockets that are already well represented in the training data, yet the
CACHE benchmarking exercise shows that hit rates on novel pockets
remain modest when the scoring function is
weak~\citep{li2024cache}. By staging VSX against VSH we allow
induced-fit effects to be captured at the decision-critical final step
without paying a flexible-receptor cost on every compound in the library,
and the fact that VSH marginally outperforms VSX on DUD confirms that
this staging is worth its computational cost. Together with the
open-source release of all three components---force field, protocol, and
platform---this moves the ultra-large-screening frontier towards
infrastructure that is also reproducible and auditable.

The campaigns also demonstrate generality across biological regimes.
KLHDC2 is a compact soluble pocket in a substrate-receptor subunit of a
CUL2/CRL2 E3 ligase, while NaV1.7 VSD4 is a membrane-embedded
state-dependent site in a voltage-gated ion channel; the two targets
differ in pocket polarity, in pocket depth, in the scoring signature
expected from bound ligands, and in the assay modality used to
validate hits. The platform recovered novel chemotypes in both settings,
in each case within seven days of wall-clock compute, and in the NaV1.7
case the selected compounds had Tanimoto similarity below 0.33 to known
NaV1.7 inhibitors, meaning that the hits were not variations on known
aryl-sulfonamide scaffolds but structurally distinct chemotypes. The
state-dependent profile of Z8739902234---inactivated-state IC$_{50}$ of
1.32~\textmu M versus resting-state IC$_{50}$ of 13.1~\textmu M---and its
modest but measurable selectivity over NaV1.5 (5.14~\textmu M) and hERG
(6.27~\textmu M) further show that the platform can produce
pharmacologically useful starting points, not only compounds that score
well in silico.

Several limitations of the current study point to concrete improvements.
The NaV1.7 campaign achieved only modest selectivity over NaV1.5 and
hERG, consistent with the well-known difficulty of the VSD4 pocket, and
would benefit from a second-round screen or a medicinal-chemistry
follow-up using the newly discovered chemotypes as starting points.
The KLHDC2 hits are in the single-digit-micromolar range, which we
anticipate could be improved to the nanomolar regime through further
compound optimisation guided by the co-crystal structure, although this
remains to be demonstrated. Finally, both campaigns relied on a known or
previously modelled receptor conformation; the platform does not yet
explicitly account for large backbone rearrangements that go beyond
sidechain flexibility, and extending VSH to include principled
backbone sampling is an open problem.

Several promising directions for future work follow from these
limitations. GPU acceleration of Rosetta docking would reduce the
wall-clock cost of both VSX and VSH and make screens of this scale
routine rather than week-long efforts on shared infrastructure.
Generative AI models for initial pose sampling could replace parts of the
VSX sampling schedule while preserving a physics-based scoring function
for the final rank-ordering. The active-learning surrogate could be
replaced by richer models, for example models that jointly learn
protein-ligand interaction features rather than treating the ligand
alone, and a generalisable deep-learning score function could be
ensembled with RosettaGenFF-VS for improved discrimination of true
binders. A further direction is template-guided screening, in which
known non-small-molecule binders such as macrocycles or antibody loops
are used as geometric priors to focus small-molecule sampling on the
relevant subpocket. The recent surge of artificial-intelligence methods
across structural biology and drug discovery suggests that integrating
these techniques with an accurate, open, flexible-receptor physics-based
core of the kind we describe here will be a productive path forward.

Finally, the broader significance of this work lies in the combination of
accuracy, openness, and scale. Ultra-large virtual screens against
multi-billion-compound libraries are no longer the exclusive province of
proprietary pipelines or purely rigid-receptor workflows. With an open
force field tuned for virtual screening, a paired VSX/VSH docking
protocol, and an active-learning platform that scales across a
3000-CPU HPC cluster, a single laboratory can now complete a
billion-compound campaign against a challenging target in under seven
days and recover experimentally validated novel hits. This makes
structure-based virtual screening at billion-compound scale available to
a much broader community and accelerates the discovery of chemotypes for
targets, such as KLHDC2 and NaV1.7, where traditional approaches have
struggled.

\bibliographystyle{unsrtnat}
\bibliography{references}

\end{document}
